% Integration fragment only. Source: low_collar_tail_paper_20261103/paper.tex,
% SHA256 e446c83bd23753919ecbe1dd1403a08858bc08e42259375135e0bdd02746026a.
% All notation and theorem environments are supplied by the v3 preamble.
\section{Bounded-window prime-tail compression and finite radical certificates}
\label{vfourB:sec}\label{vfourB:section}

This section integrates the bounded-window construction of
\cite{vfourLowWindow}. Its purpose is a quantitative finite-rank
approximation of the actual complete prime tail, followed by an exact
finite weighted-Gram certificate format. It does not assign a sign to the
low window. The separate all-Schatten result in
Theorem~\ref{vfourD:schatten} gives the stronger negative-spectral-count
input; its proof is not repeated here.

\subsection{The centered operator, conventions, and complete tail}
\label{vfourB:setting}

We use the actual spaces, kernel, and forms of
\eqref{a:eq:weights}--\eqref{a:eq:X}. Throughout this section inner products
are \emph{linear in the first argument}, as in the common notation:
\[
 \ip{u}{v}_{\Hh}=\int_{\R}u(x)\overline{v(x)}\,\frac{dx}{W(x)},
 \qquad \ip{a}{b}_{\C^n}=\sum_{k=1}^n a_k\overline{b_k}.
\]
In particular, if $Z:\C^n\to\Hh_c$ has columns $z_k$, then
$(Z^*u)_k=\ip{u}{z_k}_{\Hh}$ and
$(Z^*Z)_{ab}=\ip{z_b}{z_a}_{\Hh}$. All adjoints and Gram matrices below
use this convention.

Let $\mathsf C u=u-2m(u)\kappa$ be the weighted orthogonal projection
in~\eqref{a:eq:centering}. Write $L_c$ for the self-adjoint representative
of $q_0$ on $X_c\subset\Hh_c$, and set
$D_c=\mathsf C D_b|_{\Hh_c}$, where $D_b$ is the complete bounded
correction of~\eqref{eq:Db-physical}. Thus
\begin{equation}\label{vfourB:centered-form}
 q_0[u]=\tg[u]-\ip{D_cu}{u}_{\Hh},\qquad
 \norm{D_c}\le d,\qquad L_c\ge-dI.
\end{equation}
Lemma~\ref{a:lem:prime} permits the explicit choice $d=10$.
The correction contains $\cg\norm u_2^2$ and every prime-power shift;
$D_c$ is not the operator called $B$ in~\eqref{eq:Db-physical}.
Lemma~\ref{a:lem:fourier} gives
\begin{equation}\label{vfourB:gamma}
 \tg[u]=\int_{\R}d_\Gamma(t)|\widehat u(t)|^2\,dt,
 \qquad d_\Gamma(t)\ge\tfrac12\log(1+4t^2).
\end{equation}
Here the Fourier transform is the ordinary unitary transform.
The logarithmic large-frequency growth also follows from the digamma
formulas in \cite[Sections 5.9 and 5.11]{DLMF}; the explicit inverse cutoff
below uses only the displayed lower bound.

By Lemma~\ref{a:lem:radical}, the linearly independent vectors
$r_j=\kappa^{(2j)}-4^{-j}\kappa$ are full-domain radicals. Put
\begin{equation}\label{vfourB:known-radical}
 \Nn=\overline{\operatorname{span}\{r_j:j\ge1\}}^{\Hh},
 \qquad \Kk=\Hh_c\ominus\Nn,
 \qquad P_0=\one_{\{0\}}(L_c).
\end{equation}
The representation theorem and closedness imply
$\Nn\subseteq\ker L_c$ and $L_cP_\Nn=0$. We do not identify
$P_\Nn$ with $P_0$.

For an integer $M\ge1$, define the actual infinite prime tail by
\begin{equation}\label{vfourB:tail-form}
 p_{>M}[u]=2\operatorname{Re}\sum_{n>M}
 a_n\int_{\R}u(x+\log n)\overline{u(x)}\,dx
 =\ip{D_Mu}{u}_{\Hh},\qquad u\in\Hh_c.
\end{equation}
Its bounded self-adjoint representative $D_M$ is the centered raw weighted
row, with both orientations paired. We record the complete-tail caps used
in \cite{vfourLowWindow}:
\begin{equation}\label{vfourB:tail-caps}
 \begin{gathered}
 W(x)\le40e^{4|x|-\pi e^{2|x|}},\qquad
 W_*:=\norm W_\infty\le40e^{-\pi},\\
 \norm{D_M}\le\tau_M:=80\sum_{n>M}\Lambda(n)n^{3/2}e^{-\pi n}
 \le160(M+1)^2e^{-\pi(M+1)}.
 \end{gathered}
\end{equation}
These are weighted-space norm bounds on the \emph{full} tail, not a
positivity statement.

For completeness, the constant $40$ follows from the envelope in
Lemma~\ref{a:lem:theta}. In its defining series, the terms
$j^4e^{-\pi(j^2-1)}$ for $j\ge2$ have successive ratio at most
$(3/2)^4e^{-5\pi}<1/1000$, and the first is less than $1/500$.
Their sum is therefore less than $1/400$. Using $\pi<22/7$ gives
$C_\kappa<(1936/49)(401/400)<40$.
For a displacement $s\ge0$, the same weighted-shift argument as in
Lemma~\ref{a:lem:prime} gives
\[
 \sup_x\sqrt{W(x)W(x+s)}\le40e^{2s-\pi e^s}.
\]
Consequently each paired $n$th row has norm at most
$80\Lambda(n)n^{3/2}e^{-\pi n}$. Centering is contractive.
Finally $\log n\le\sqrt n$ for $n\ge1$, and the consecutive ratios of
$n^2e^{-\pi n}$ for $n\ge2$ are less than $1/2$. Summing that geometric
majorant proves the last inequality in~\eqref{vfourB:tail-caps}.
In every upper bound or parameter choice below, the explicit cap
$40e^{-\pi}$ may replace $W_*$.

\subsection{Finite-rank approximation on a bounded spectral window}
\label{vfourB:approximation}

Fix a bounded Borel set $J\subset[-R,R]$, where $R\ge0$, and write
$P_J=\one_J(L_c)$. The spectral theorem gives
$\Ran P_J\subset\dom L_c$ and
\begin{equation}\label{vfourB:window-energy}
 \tg[u]\le(R+d)\norm u_{\Hh}^2,
 \qquad \norm u_2\le\sqrt{W_*}\norm u_{\Hh}
 \qquad(u\in\Ran P_J).
\end{equation}
For integers $N\ge M$, define the finite raw sum
\[
 D_{(M,N]}^{\mathrm{raw}}u(x)
 =W(x)\sum_{M<n\le N}a_n
       \bigl(u(x+\log n)+u(x-\log n)\bigr).
\]
An individual oriented row maps ordinary $L^2(dx)$ into the ambient
weighted space $L^2(dx/W)$, since
\[
 \norm{W(\cdot)u(\cdot+s)}_{L^2(dx/W)}^2
 =\int_{\R}W(x)|u(x+s)|^2\,dx\le W_*\norm u_2^2.
\]
Pairing the orientations preserves evenness, after which $\mathsf C$ is
applied. Thus
\begin{equation}\label{vfourB:finite-shift-cap}
 \norm{\mathsf C D_{(M,N]}^{\mathrm{raw}}}
       _{\Le\to\Hh_c}
 \le C_{M,N}:=2\sqrt{W_*}\sum_{M<n\le N}\frac{\Lambda(n)}{\sqrt n}.
\end{equation}
Only this \emph{finite} shift sum is used as an ordinary-to-weighted map.
The infinite remainder is bounded in the full $\Hh$ norm by $\tau_N$;
no $L^2(dx)\to\Hh$ bound is asserted for the infinite prime tail.

Let $F_T$ be the ordinary Fourier projection onto $[-T,T]$, and set
\[
 W_B=\sup_{|x|>B}W(x),\qquad
 \nu(T)=\inf_{|t|>T}d_\Gamma(t).
\]
For $B,T>0$ and an integer $m\ge1$, define
\begin{equation}\label{vfourB:taylor-map}
 \begin{split}
 \mathcal T_{B,T,m}u(x)
 &=\one_{[-B,B]}(x)\sum_{\substack{0\le k<m\\k\ \mathrm{even}}}
       x^k\ell_k(u),\\
 \ell_k(u)&=\frac{1}{\sqrt{2\pi}\,k!}
            \int_{-T}^T(it)^k\widehat u(t)\,dt.
 \end{split}
\end{equation}
The output is used only in ordinary $L^2$, and the map has rank at most
$\lceil m/2\rceil$. Define
\begin{equation}\label{vfourB:error-parameters}
 \zeta_{B,T,m}=\sqrt{\frac{2BT}{\pi}}\frac{(BT)^m}{m!},\qquad
 b_{R;B,T,m}=\sqrt{W_B}
       +\sqrt{\frac{R+d}{\nu(T)}}+\sqrt{W_*}\,\zeta_{B,T,m}.
\end{equation}

\begin{theorem}[Actual bounded-window tail approximation]
\label{vfourB:approximation-theorem}
Assume $\nu(T)>0$, and put
\[
 A=\mathsf C D_{(M,N]}^{\mathrm{raw}}
                   \mathcal T_{B,T,m}|_{\Hh_c},\qquad
 E=\tau_N+C_{M,N}b_{R;B,T,m}.
\]
Then $A$ is bounded on $\Hh_c$, has rank at most $\lceil m/2\rceil$, and
\begin{equation}\label{vfourB:row-error}
 \norm{(D_M-A)P_J}\le E.
\end{equation}
For the self-adjoint finite-rank operator $S=(A+A^*)/2$, one has the
\emph{compressed} estimate
\begin{equation}\label{vfourB:compressed-error}
 \norm{P_J(D_M-S)P_J}\le E.
\end{equation}
No corresponding estimate for the symmetrized row $(D_M-S)P_J$ is claimed.
\end{theorem}
\begin{proof}
For $u\in\Ran P_J$, the spatial and Fourier tails satisfy
\[
 \norm{\one_{\{|x|>B\}}u}_2\le\sqrt{W_B}\norm u_{\Hh},\qquad
 \norm{(I-F_T)u}_2
 \le\sqrt{\frac{R+d}{\nu(T)}}\norm u_{\Hh}.
\]
The integral Taylor remainder for real $tx$ is bounded by
$|tx|^m/m!$. The resulting kernel error on
$[-B,B]\times[-T,T]$, including the factor $(2\pi)^{-1/2}$, has
Hilbert--Schmidt norm at most $\zeta_{B,T,m}$. The odd terms vanish on even
inputs. Decomposing
\[
 u-\mathcal T_{B,T,m}u
 =\one_{\{|x|>B\}}u+\one_{[-B,B]}(u-F_Tu)
   +\bigl(\one_{[-B,B]}F_Tu-\mathcal T_{B,T,m}u\bigr)
\]
and using~\eqref{vfourB:window-energy} therefore gives
\begin{equation}\label{vfourB:ordinary-error}
 \norm{u-\mathcal T_{B,T,m}u}_2
 \le b_{R;B,T,m}\norm u_{\Hh}.
\end{equation}
All functionals $\ell_k$ are bounded on $\Hh$ by its continuous embedding
in ordinary $L^2$, so $A$ is bounded and its rank has the stated bound.
The identity
\[
 D_M-A=D_N+\mathsf C D_{(M,N]}^{\mathrm{raw}}
                         (I-\mathcal T_{B,T,m})|_{\Hh_c}
\]
combines the complete weighted tail estimate with
\eqref{vfourB:finite-shift-cap} and~\eqref{vfourB:ordinary-error} to
prove~\eqref{vfourB:row-error}. Finally
$P_J(D_M-S)P_J$ is the self-adjoint part of $P_J(D_M-A)P_J$.
Taking a self-adjoint part cannot increase the operator norm, proving
\eqref{vfourB:compressed-error} without estimating an uncompressed
adjoint row.
\end{proof}

The columns of $A$ are centered finite prime-shift rows of the ordinary
compactly supported functions $\one_{[-B,B]}x^k$. Their construction does
not require a sharp spectral projection. The explicit schedule below
makes $E$ arbitrarily small; hence $D_MP_J$ is compact, even when $P_J$
has infinite rank.

\subsection{Explicit inverse-Gamma cutoff and the minimal rank schedule}
\label{vfourB:schedule}

Fix $0<\varepsilon\le1$ and set
\begin{equation}\label{vfourB:integer-cutoff}
 N=\max\{M,\lceil3+\log(4/\varepsilon)\rceil\}.
\end{equation}
Then $\tau_N<\varepsilon/4$. Indeed the function
$160y^2e^{-(\pi-1)y-4}$ decreases for $y\ge4$ and is less than one at
$y=4$, so~\eqref{vfourB:tail-caps} with
$y=N+1\ge4+\log(4/\varepsilon)$ suffices.
If $C_{M,N}=0$, take $A=0$, with no further parameters required.
Otherwise choose an explicit cap $C_\sharp\ge C_{M,N}>0$ and define
\begin{equation}\label{vfourB:spatial-cutoff}
 X_\varepsilon=\max\{e^2,4C_\sharp\sqrt{40}/\varepsilon\},\qquad
 Y_\varepsilon=\frac4\pi\log X_\varepsilon,\qquad
 B=\tfrac12\log Y_\varepsilon.
\end{equation}
The inverse of the logarithmic Gamma lower bound gives the explicit
frequency cutoff
\begin{equation}\label{vfourB:inverse-gamma}
 H_\varepsilon=
 \max\left\{1,\frac{16C_\sharp^2(R+d)}{\varepsilon^2}\right\},\qquad
 T=\frac12\sqrt{e^{2H_\varepsilon}-1}.
\end{equation}
In particular $\nu(T)\ge H_\varepsilon>0$, including when $R+d=0$.
Choose $m$ to be the \emph{smallest positive integer} satisfying
\begin{equation}\label{vfourB:minimal-m}
 \begin{split}
 m&\ge4eBT,\\
 m&\ge\max\left\{0,
   \log_4\!\left(
       \frac{4C_\sharp\sqrt{W_*}\sqrt{2BT/\pi}}{\varepsilon}
           \right)\right\}.
 \end{split}
\end{equation}
Equivalently, $m$ is the maximum of $1$ and the ceilings of the two
right-hand sides.

\begin{proposition}[Complete error budget and rank cost]
\label{vfourB:rank-schedule}
The preceding choices give $E\le\varepsilon$ in
Theorem~\ref{vfourB:approximation-theorem}. With the specified minimal $m$
and the explicit cap
\begin{equation}\label{vfourB:chebyshev-cap}
 C_\sharp=2\sqrt{W_*}\,[4(\log2)\sqrt N+5],
\end{equation}
for fixed $M,R,d,W_*$ the construction satisfies
\begin{equation}\label{vfourB:rank-bound}
 \rank A\le\lceil m/2\rceil
 \le\exp\!\left(O_{M,R,d,W_*}
       \bigl(\varepsilon^{-2}[1+\log(1/\varepsilon)]\bigr)\right).
\end{equation}
\end{proposition}
\begin{proof}
For $X_\varepsilon\ge e^2$ one has
$Y_\varepsilon\le X_\varepsilon$. The decreasing theta envelope gives
\[
 C_\sharp\sqrt{W_B}
 \le C_\sharp\sqrt{40}\,Y_\varepsilon
                  e^{-\pi Y_\varepsilon/2}
 \le\frac{C_\sharp\sqrt{40}}{X_\varepsilon}
 \le\frac{\varepsilon}{4}.
\]
Equation~\eqref{vfourB:inverse-gamma} gives
$C_\sharp\sqrt{(R+d)/\nu(T)}\le\varepsilon/4$.
The elementary factorial estimate $m!\ge(m/e)^m$ and
\eqref{vfourB:minimal-m} give
\[
 \zeta_{B,T,m}\le\sqrt{2BT/\pi}\,4^{-m},\qquad
 C_\sharp\sqrt{W_*}\,\zeta_{B,T,m}\le\varepsilon/4.
\]
Together with the discarded complete tail, these are the four parts of
$E\le\varepsilon$.

To check~\eqref{vfourB:chebyshev-cap}, write
$\psi_{\mathrm C}(x)=\sum_{n\le x}\Lambda(n)$.
The central-binomial estimate yields
$\psi_{\mathrm C}(n)\le\psi_{\mathrm C}(\lceil n/2\rceil)+n\log2$
for integers $n\ge2$. Iterating down to $1$ gives, for real $x\ge1$,
\[
 \psi_{\mathrm C}(x)\le2(\log2)x+\log x+\log2.
\]
Partial summation then gives
\[
 \sum_{n\le N}\frac{\Lambda(n)}{\sqrt n}
 =\frac{\psi_{\mathrm C}(N)}{\sqrt N}
   +\frac12\int_1^N\frac{\psi_{\mathrm C}(t)}{t^{3/2}}\,dt
 \le4(\log2)\sqrt N+5,
\]
which proves the cap. As $\varepsilon\downarrow0$,
$N=O_M(1+\log(1/\varepsilon))$ and
$C_\sharp^2=O_{M,W_*}(1+\log(1/\varepsilon))$.
Hence
$H_\varepsilon=O_{M,R,d,W_*}
(\varepsilon^{-2}[1+\log(1/\varepsilon)])$,
$T\le e^{H_\varepsilon}/2$, and $B$ grows at most logarithmically in
$1+\log(1/\varepsilon)$. The logarithmic second requirement
in~\eqref{vfourB:minimal-m} is $O(H_\varepsilon+1+\log(1/\varepsilon))$.
Taking the \emph{smallest} integer satisfying both requirements now gives
\eqref{vfourB:rank-bound}; enlarging constants covers the remaining compact
range of $\varepsilon$.
\end{proof}

An arbitrarily inflated cap $C_\sharp$, or an arbitrarily enlarged
admissible $m$, is not covered by the stated rank estimate. The construction
is not claimed to be optimal or practical. In particular, the exponential
$e^{2H_\varepsilon}$ in~\eqref{vfourB:inverse-gamma} cannot be replaced by
$e^{H_\varepsilon}$ when inverting the displayed Gamma bound.

\subsection{Continuous filters and exact finite constrained Grams}
\label{vfourB:gram-section}

Let $I\subset J$, and choose a real continuous function $0\le\phi\le1$
that equals one on $I$ and has support in $J$, whenever such a filter is
available. Put
\begin{equation}\label{vfourB:filtered-operator}
 F=\phi(L_c),\qquad S_\phi=\tfrac12 F(A+A^*)F.
\end{equation}
Since $FP_J=F=P_JF$, the compressed estimate implies the global bound
\begin{equation}\label{vfourB:filtered-error}
 \norm{FD_MF-S_\phi}\le E.
\end{equation}
Thus $F$ is the functional calculus of the \emph{complete} centered
operator, not of a Gamma-only substitute.

Let $r=\lceil m/2\rceil$ and index its columns by the even integers
$0\le k<m$. Set
\[
 y_k=\mathsf C D_{(M,N]}^{\mathrm{raw}}
                     (\one_{[-B,B]}x^k),\qquad
 g_k=\mathsf C(Wb_k),\qquad
 \widehat b_k(t)=\frac{(-it)^k}{\sqrt{2\pi}\,k!}\one_{[-T,T]}(t).
\]
Here $b_k\in\Le$, $Wb_k\in\Hh$, and Plancherel gives
$\ell_k(u)=\ip{u}{g_k}_{\Hh}$ for $u\in\Hh_c$ under the declared
linear-first convention. Consequently, with column maps
$Y,G:\C^r\to\Hh_c$, one has $A=YG^*$. Define
\begin{equation}\label{vfourB:factorization}
 Z=(FY,FG):\C^{2r}\to\Hh_c,\qquad
 H_0=\frac12\begin{pmatrix}0&I_r\\ I_r&0\end{pmatrix},\qquad
 S_\phi=ZH_0Z^*.
\end{equation}
If $A=0$ was chosen above, take empty column maps instead.

Let $R_j:\C^j\to\Hh_c$ have columns $r_1,\ldots,r_j$, and let $P_j$
be the true $\Hh$ projection onto their span. Put
\begin{equation}\label{vfourB:constraint-gram}
 \begin{gathered}
 H_j=R_j^*R_j>0,\qquad B_j=R_j^*Z,\qquad
 P_j=R_jH_j^{-1}R_j^*,\qquad Q_j=I-P_j,\\
 G_j=Z^*Q_jZ=Z^*Z-B_j^*H_j^{-1}B_j\ge0.
 \end{gathered}
\end{equation}
For $j=0$ the constraint term is absent. The entries of these matrices
are the actual reciprocal-theta weighted inner products. In particular,
$G_j$ is not a Euclidean Gram of unweighted coefficients.

\begin{theorem}[Finite known-radical tail certificates]
\label{vfourB:gram-certificate}
Define
\[
 K_j=G_j^{1/2}H_0G_j^{1/2},\qquad
 \alpha_j^+=\max\{0,\lambda_{\max}(K_j)\},\qquad
 \alpha_j^-=\max\{0,-\lambda_{\min}(K_j)\},
\]
with both extrema defined to be zero for an empty matrix. For every
$u\in\Kk$ satisfying $Fu=u$,
\begin{equation}\label{vfourB:two-sided-certificate}
 -(E+\alpha_j^-)\norm u_{\Hh}^2
 \le p_{>M}[u]\le(E+\alpha_j^+)\norm u_{\Hh}^2.
\end{equation}
Singular $G_j$ are allowed. The two sequences $\alpha_j^\pm$ are
nonincreasing, and
\begin{equation}\label{vfourB:constraint-limit}
 Q_jS_\phi Q_j\longrightarrow P_\Kk S_\phi P_\Kk
 \quad\hbox{in operator norm as }j\to\infty.
\end{equation}
No effective rate of this convergence is asserted.
\end{theorem}
\begin{proof}
The identity for $P_j$ follows from the positive definiteness of the
true radical Gram, using the linear independence in
Lemma~\ref{a:lem:radical}. It proves~\eqref{vfourB:constraint-gram}.
The polar decomposition $Q_jZ=V_jG_j^{1/2}$ gives
\[
 Q_jS_\phi Q_j=V_jK_jV_j^*.
\]
The partial isometry $V_j$ is isometric on the support of $G_j$, and
$K_j$ vanishes on its kernel. Thus the two operators have the same
nonzero eigenvalues, with multiplicities, and their quadratic forms are
bounded between $-\alpha_j^-I$ and $\alpha_j^+I$ after adjoining the
possible zero subspaces. For $u\in\Kk$, $Q_ju=u$; for $Fu=u$,
\eqref{vfourB:filtered-error} bounds the difference between
$p_{>M}[u]$ and $\ip{S_\phi u}{u}_{\Hh}$ by $E\norm u_{\Hh}^2$.
This proves~\eqref{vfourB:two-sided-certificate}.

The quantities $\alpha_j^\pm$ are the nonnegative Rayleigh extrema of
$\pm S_\phi$ on the nested subspaces $\Ran Q_j$, so they cannot increase
with $j$. The projections $P_j$ increase strongly to $P_\Nn$.
Since $S_\phi$ has finite rank, strong convergence of these bounded
projections implies the norm convergence in
\eqref{vfourB:constraint-limit}.
\end{proof}

These are safe \emph{absolute} upper bounds for both tail signs, not a
certified sign. A numerical certificate must additionally enclose the
actual filtered columns, the weighted Gram entries, the radical inverse,
and the finite matrix extrema, and must add all operator errors to $E$.
No such numerical evaluation is supplied here.

\subsection{A resolvent implementation of the filter}
\label{vfourB:resolvent-filter}

The sharp projection $P_J$ is needed only to prove the error estimate.
The filter can instead be approximated from actual resolvents. Suppose
$\phi$ is Lipschitz and supported in $[-R,R]$. For $s>d$, set
\[
 R_s=(L_c+s)^{-1},\qquad a=(s-d)^{-1},\qquad
 h(r)=\begin{cases}\phi(1/r-s),&r>0,\\0,&r=0.\end{cases}
\]
On $[0,a]$, the function $h$ is Lipschitz with constant at most
$\operatorname{Lip}(\phi)(s+R)^2$. Indeed it vanishes for
$r<1/(s+R)$, and on the remaining interval the reciprocal map has
Lipschitz constant at most $(s+R)^2$; continuity handles the joining
point. Its degree-$q$ Bernstein polynomial, scaled to $[0,a]$, is
\[
 p_q(r)=\sum_{k=0}^q h(ak/q)\binom qk(r/a)^k(1-r/a)^{q-k}.
\]
For $q\ge1$, the binomial variance bound and the spectral theorem imply
\begin{equation}\label{vfourB:bernstein-error}
 \norm{p_q(R_s)-F}
 \le\frac{\operatorname{Lip}(\phi)(s+R)^2}
                 {2(s-d)\sqrt q}.
\end{equation}
To see this, at $r=a\theta$ write $p_q(r)=\mathbb E h(aK/q)$ with
$K$ binomial of parameters $q,\theta$. The Lipschitz error is at most
$a\operatorname{Lip}(h)\sqrt{\theta(1-\theta)/q}$, uniformly bounded
by the right-hand side. Since $h(0)=0$, $p_q$ has zero constant term.

Actual resolvent columns still require verification. If
$v\in\dom L_c$ and $f\in\Hh_c$, then
\begin{equation}\label{vfourB:resolvent-residual}
 \norm{v-R_sf}_{\Hh}
 \le\frac{\norm{(L_c+s)v-f}_{\Hh}}{s-d}.
\end{equation}
If a bounded approximation $\widetilde F$ satisfies
$\norm{\widetilde F-F}\le\eta$, the self-adjoint replacement
$\widetilde F S\widetilde F^*$, where $S=(A+A^*)/2$, obeys
\begin{equation}\label{vfourB:filter-replacement}
 \norm{\widetilde F S\widetilde F^*-FSF}
 \le\eta(2+\eta)\norm S
 \le\eta(2+\eta)\norm A.
\end{equation}
Indeed expand the difference into two terms and use
$\norm F\le1$, $\norm{\widetilde F}\le1+\eta$.
For a self-adjoint $\widetilde F$, such as the exact real polynomial
$p_q(R_s)$, the adjoint on the right factor is redundant. These estimates
specify what must be certified; they do not supply computed resolvent or
filter columns.

\subsection{Kernel-free collars and the remaining zero-fibre distinction}
\label{vfourB:collars}

\begin{proposition}[Two different collar limits]
\label{vfourB:collar-limits}
Fix $M$ and a bounded Borel $J$ containing a neighborhood of zero. Then
\begin{equation}\label{vfourB:kernel-free-limit}
 E_\delta=\one_{[-\delta,\delta]\setminus\{0\}}(L_c)
 \quad\Longrightarrow\quad
 \norm{D_ME_\delta}\longrightarrow0
 \quad(\delta\downarrow0).
\end{equation}
In contrast, for continuous filters $0\le\phi_\delta\le1$ equal to one
on $[-\delta,\delta]$ and supported in $[-2\delta,2\delta]$,
\begin{equation}\label{vfourB:known-radical-limit}
 P_\Kk\phi_\delta(L_c)D_M\phi_\delta(L_c)P_\Kk
 \longrightarrow (P_0-P_\Nn)D_M(P_0-P_\Nn)
 \quad\hbox{in operator norm}.
\end{equation}
The second limit need not vanish unless additional information about
the residual zero fibre or its tail compression is available.
\end{proposition}
\begin{proof}
The spectral theorem gives $E_\delta\to0$ strongly, with
$P_JE_\delta=E_\delta$ for sufficiently small $\delta$.
Theorem~\ref{vfourB:approximation-theorem} and
Proposition~\ref{vfourB:rank-schedule} imply that $D_MP_J$ is compact.
A compact operator multiplied by this strongly null family of
self-adjoint contractions converges to zero in norm, proving
\eqref{vfourB:kernel-free-limit}.

Also $F_\delta=\phi_\delta(L_c)\to P_0$ strongly by dominated convergence
in the spectral theorem, and $F_\delta=P_JF_\delta$ for small $\delta$.
Since $P_JD_MP_J$ is compact,
\[
 F_\delta(P_JD_MP_J)F_\delta\longrightarrow P_0D_MP_0
 \quad\hbox{in norm}.
\]
The known-radical subspace lies in $\ker L_c$, so its projection commutes
with the functional calculus and $P_\Kk P_0=P_0-P_\Nn$.
Compressing the last limit by $P_\Kk$ proves
\eqref{vfourB:known-radical-limit}.
\end{proof}

The kernel-free conclusion is only an absolute $o(1)$ statement.
Neither compactness nor this construction supplies $O(\delta)$ or
$o(\delta)$ control. In particular, choosing a target error
$\varepsilon=c\delta$ in the explicit schedule does not prove that the
finite matrix contribution in~\eqref{vfourB:two-sided-certificate} is
$O(\delta)$. The stronger all-Schatten result of
Theorem~\ref{vfourD:schatten} addresses negative spectral counting, not
the missing relative low-collar tail sign. No full-kernel identification,
whole-space sign conclusion, or proof or disproof of the Riemann
hypothesis follows from this section.
